%==============================================================================
%  A Practical Guide to AXI4-Stream and AXI4-Stream Video Interfaces
%
%  @author : Yosel Balibrea Lastre
%  @date   : 2026-08-28
%  @brief  : Skeleton / outline for the AXI4-Stream video tutorial.
%
%  Build:  pdflatex report && pdflatex report      (twice for the ToC/refs)
%==============================================================================
\documentclass[11pt,a4paper,oneside]{report}

%------------------------------------------------------------------ encoding --
\usepackage[T1]{fontenc}
\usepackage[utf8]{inputenc}
\usepackage[english]{babel}
\usepackage{lmodern}
\usepackage{microtype}

%------------------------------------------------------------------- layout ---
\usepackage[a4paper,margin=2.5cm,headheight=14pt]{geometry}
\usepackage{fancyhdr}
\usepackage{parskip}
\usepackage{enumitem}
\usepackage{float}

%------------------------------------------------------------------ graphics --
\usepackage{graphicx}
\usepackage{xcolor}

%-------------------------------------------------------------------- colours --
\definecolor{NavyBlue}{HTML}{1F4E79}
\definecolor{CodeBg}{HTML}{F7F7F7}
\definecolor{CodeKw}{HTML}{0000B0}
\definecolor{CodeCm}{HTML}{4E8B4E}
\definecolor{CodeStr}{HTML}{A31515}
\definecolor{CodeRule}{HTML}{D0D0D0}
\usepackage{tikz}
\usetikzlibrary{arrows.meta,positioning,calc,fit,backgrounds,shapes.multipart}

% tikz-timing is ideal for handshake waveforms but is not in every TeX install.
\newif\ifhastimingpkg
\IfFileExists{tikz-timing.sty}{\usepackage{tikz-timing}\hastimingpkgtrue}
                              {\hastimingpkgfalse}

%-------------------------------------------------------------------- tables --
\usepackage{booktabs}
\usepackage{longtable}
\usepackage{array}
\usepackage{multirow}
\usepackage{tabularx}
\newcolumntype{L}[1]{>{\raggedright\arraybackslash}p{#1}}

%---------------------------------------------------------------- math/units --
\usepackage{amsmath}
\usepackage{amssymb}
\usepackage{siunitx}
\sisetup{per-mode=symbol,detect-weight=true}

%-------------------------------------------------------------------- floats --
\usepackage{caption}
\usepackage{subcaption}
\captionsetup{font=small,labelfont=bf}

%--------------------------------------------------------------------- links --
\usepackage[colorlinks=true,
            linkcolor=NavyBlue,
            citecolor=NavyBlue,
            urlcolor=NavyBlue,
            pdfauthor={Yosel Balibrea Lastre},
            pdftitle={A simple AXI4-Stream Video Pattern Generator}]{hyperref}
\usepackage{cleveref}

%----------------------------------------------------------------- listings ---
\usepackage{listings}

\lstdefinestyle{verilog}{
  language        = Verilog,
  basicstyle      = \ttfamily\footnotesize,
  keywordstyle    = \color{CodeKw}\bfseries,
  commentstyle    = \color{CodeCm}\itshape,
  stringstyle     = \color{CodeStr},
  backgroundcolor = \color{CodeBg},
  frame           = single,
  rulecolor       = \color{CodeRule},
  numbers         = left,
  numberstyle     = \tiny\color{gray},
  numbersep       = 8pt,
  breaklines      = true,
  showstringspaces= false,
  tabsize         = 4,
  captionpos      = b,
  morekeywords    = {logic,always_ff,always_comb,typedef,interface,modport,
                     package,import,endpackage,endinterface,bit,byte},
}
\lstset{style=verilog}

% Inline signal names:  \sig{TVALID}
\newcommand{\sig}[1]{\texttt{#1}}
% Inline module/file names: \file{axis_skid_buffer.v}
\newcommand{\file}[1]{\texttt{#1}}

%------------------------------------------------------------------- boxes ----
\usepackage[most]{tcolorbox}

\newtcolorbox{rulebox}[1][]{colback=NavyBlue!4,colframe=NavyBlue!60,
  fonttitle=\bfseries,title=Protocol rule,#1}
\newtcolorbox{pitfall}[1][]{colback=red!3,colframe=red!50!black,
  fonttitle=\bfseries,title=Common pitfall,#1}
\newtcolorbox{notebox}[1][]{colback=black!3,colframe=black!40,
  fonttitle=\bfseries,title=Note,#1}

%------------------------------------------------------------------ headers ---
\pagestyle{fancy}
\fancyhf{}
\fancyhead[L]{\small\nouppercase{\leftmark}}
\fancyhead[R]{\small\thepage}
\renewcommand{\headrulewidth}{0.4pt}

%---------------------------------------------------------------- TODO macro --
% Set \showtodofalse before the final build to hide all placeholders.
\newif\ifshowtodo\showtodotrue
\newcommand{\TODO}[1]{\ifshowtodo\textcolor{red!70!black}{\small\textbf{[TODO]}~#1}\fi}

\setcounter{tocdepth}{2}
\setcounter{secnumdepth}{3}

%==============================================================================
\begin{document}
%==============================================================================

% Front matter uses roman numerals so PDF page anchors stay unique.
\pagenumbering{roman}

\begin{titlepage}
  \centering
  \vspace*{3cm}
  {\Huge\bfseries A simple AXI4-Stream Video Pattern Generator\par}
  \vspace{1cm}
  {\Large Interfaces, protocol rules and RTL design patterns\par}
  \vspace{2.5cm}
  {\large Yosel Balibrea Lastre\par}
  \vspace{0.5cm}
  {\large \today\par}
  \vfill
  {\small Revision 0.1 --- draft\par}
\end{titlepage}
\setcounter{page}{2}   % report.cls resets to 1 here; keep anchors unique

\tableofcontents
\listoffigures
\listoftables

%==============================================================================
\chapter*{About this guide}
\addcontentsline{toc}{chapter}{About this guide}
\markboth{About this guide}{}
%==============================================================================


\section*{Accompanying source code}
\TODO{Point at the repository layout, e.g.\ \file{m\_axis.v}, \file{s\_axis.v},
testbenches and simulation scripts.}

\cleardoublepage
\pagenumbering{arabic}

%==============================================================================
\part{AXI4-Stream Fundamentals}
%==============================================================================

%------------------------------------------------------------------------------
\chapter{Introduction to Streaming Interfaces}
%------------------------------------------------------------------------------

\section{Why a streaming interface?}
\TODO{Memory-mapped vs.\ streaming data movement; when each is appropriate.}

\section{Where AXI4-Stream sits in the AMBA family}
\TODO{AXI4 / AXI4-Lite / AXI4-Stream comparison table.}

\begin{table}[H]
  \centering
  \caption{Comparison of the AMBA AXI interface flavours.}
  \label{tab:axi-flavours}
  \begin{tabularx}{\textwidth}{l X X X}
    \toprule
    & \textbf{AXI4} & \textbf{AXI4-Lite} & \textbf{AXI4-Stream} \\
    \midrule
    Addressing      & \TODO{} & \TODO{} & \TODO{} \\
    Channels        & \TODO{} & \TODO{} & \TODO{} \\
    Bursts          & \TODO{} & \TODO{} & \TODO{} \\
    Typical use     & \TODO{} & \TODO{} & \TODO{} \\
    \bottomrule
  \end{tabularx}
\end{table}

\section{Terminology: transfer, packet, frame, stream}
\TODO{Define each level of the hierarchy and how they map to signals.}

\section{Roadmap of the document}
\TODO{One paragraph per part.}

%------------------------------------------------------------------------------


\end{document}
